\subsection{The Krull-Akizuki Theorem}

LEMMA 1. Let A be a Noetherian domain in which every non-zero prime ideal is maximal and M a finitely generated torsion A-module. Then the length long \(_{A}\) (M) of M is finite.

As M is a torsion module, every prime ideal associated with M is \(\neq(0)\) and therefore maximal. The lemma then follows from Chapter IV, §2, no. 5, Proposition 7.

LEMMA 2. Let A be a ring, T an A-module and \((\mathrm{T}_{\iota})\) a right directed family of submodules of T with union T. Then \(\operatorname{long}_{\mathbf{A}}(\mathbf{T}) = \sup (\operatorname{long}_{\mathbf{A}}(\mathbf{T}_{\iota}))\) .

\(\text{long}_{\mathbf{A}}(\mathbf{T}_{\iota}) \leqslant \text{long}_{\mathbf{A}}(\mathbf{T})\) for all \(\iota\) . The lemma is obvious if no integer exceeds the \(\text{long}_{\mathbf{A}}(\mathbf{T}_{\iota})\) , both sides then being infinite. Otherwise, let \(\iota_{0}\) be an index for which \(\text{long}_{\mathbf{A}}(\mathbf{T}_{\iota})\) takes its greatest value; then \(\mathbf{T}_{\iota_{0}} = \mathbf{T}\) since the family \((\mathbf{T}_{\iota})\) is directed; whence our assertion in this case.

Remark. This proof does not assume that A is commutative.

LEMMA 3. Let A be a Noetherian domain such that every non-zero prime ideal of A is maximal, M a torsion-free A-module of finite rank r and a non-zero element of A. Then A/Aa is an A-module of finite length and:

\[
\operatorname{long} _ {\mathbf {A}} (\mathrm{M} / a \mathrm{M}) \leqslant r. \operatorname{long} _ {\mathbf {A}} (\mathrm{A} / \mathrm{A} a).\tag{6}
\]

Lemma 1 shows that \(\mathrm{long}_{\mathbf{A}}(\mathbf{A} / \mathbf{A}a)\) is finite. We show (6) first in the case where \(\mathbf{M}\) is finitely generated. As \(\mathbf{M}\) is torsion-free and of rank \(r\) , there exists a submodule \(\mathbf{L}\) of \(\mathbf{M}\) which is isomorphic to \(\mathbf{A}'\) and such that \(\mathbf{Q} = \mathbf{M} / \mathbf{L}\) is a finitely generated torsion A-module and hence of finite length (Lemma 1). For every integer \(n \geqslant 1\) , the kernel of the canonical surjection \(\mathbf{M} / a^n\mathbf{M} \to \mathbf{Q} / a^n\mathbf{Q}\) is equal to \((\mathbf{L} + a^n\mathbf{M}) / a^n\mathbf{M}\) and isomorphic to \(\mathbf{L} / (a^n\mathbf{M} \cap \mathbf{L})\) ; as

\[
a ^ {n} \mathbf {L} \subset a ^ {n} \mathbf {M} \cap \mathbf {L},
\]

therefore

\[
\begin{array}{r l} \operatorname{long} _ {\mathbf {A}} (\mathrm{M} / a ^ {n} \mathrm{M}) & \leqslant \\ \operatorname{long} _ {\mathbf {A}} (\mathrm{L} / a ^ {n} \mathrm{L}) + \operatorname{long} _ {\mathbf {A}} (\mathrm{Q} / a ^ {n} \mathrm{Q}) & \leqslant \operatorname{long} _ {\mathbf {A}} (\mathrm{L} / a ^ {n} \mathrm{L}) + \operatorname{long} _ {\mathbf {A}} (\mathrm{Q}). \end{array} \tag {7}
\]

Now, since M is torsion-free, multiplication by a defines an isomorphism of

M/aM onto \(aA/a^{2}M\) ; similarly for L; whence, by induction on n, the formulae:

\[
\operatorname{long} _ {\mathbf {A}} (\mathbf {M} / a ^ {n} \mathbf {M}) = n. \operatorname{long} _ {\mathbf {A}} (\mathbf {M} / a \mathbf {M}).\tag{8}
\]

\[
\operatorname{long} _ {\mathbf {A}} (\mathbf {L} / a ^ {n} \mathbf {L}) = n. \operatorname{long} _ {\mathbf {A}} (\mathbf {L} / a \mathbf {L}).
\]

Taking account of (7) we deduce:

\[
\operatorname{long} _ {\mathbf {A}} (\mathrm{M} / a \mathrm{M}) \leqslant \operatorname{long} _ {\mathbf {A}} (\mathrm{L} / a \mathrm{L}) + n ^ {- 1} \operatorname{long} _ {\mathbf {A}} (\mathrm{Q})\tag{9}
\]

for all \(n > 0\) ; as \(L\) is isomorphic to \(A^r\) , \(\text{long}_A(L/aL) = r \text{ long}_A(A/Aa)\) ; whence (6) by letting \(n\) tend to infinity in (9).

We now pass to the general case. Let \((\mathbf{M}_{\iota})\) be the family of finitely generated submodules of M. The module T = M/aM is the union of the submodules \(\mathrm{T}_{\iota} = (\mathrm{M}_{\iota} + a\mathrm{M})/a\mathrm{M} = \mathrm{M}_{\iota}/(\mathrm{M}_{\iota} \cap a\mathrm{M})\) . Now, \(T_{\iota}\) is isomorphic to a quotient of \(M_{\iota}/aM_{\iota}\) and hence

\[
\operatorname{long} _ {\mathbf {A}} \left(\mathrm{T} _ {t}\right) \leqslant r \operatorname{long} _ {\mathbf {A}} (\mathrm{A} / \mathrm{A} a)
\]

by what we have just proved. Whence

\[
\operatorname{long} _ {\mathbf {A}} (\mathrm{T}) \leqslant r \operatorname{long} _ {\mathbf {A}} (\mathrm{A/Aa})
\]

by Lemma 2.

PROPOSITION 5 (Krull-Akizuki). Let A be a Noetherian domain each of whose nonzero prime ideals is maximal, K its field of fractions, L a finite extension of K and B a subring of L containing A. Then B is Noetherian and every non-zero prime ideal of B is maximal. Moreover, for every ideal b ≠ (0) of B, B/b is a finitely generated A-module.

Let b be a non-zero ideal of B. We shall show that B/b is an A-module of finite length (hence, a fortiori, a B-module of finite length) and that b is a finitely generated B-module.

A non-zero element y of b satisfies an equation of the form:

\[
a _ {r} y ^ {r} + a _ {r - 1} y ^ {r - 1} + \dots + a _ {0} = 0 \quad (a _ {1} \in A, a _ {0} \neq 0).
\]

This equation shows that \(a_{0} \in By \subset b\) . Applying Lemma 3 to M = B, it is seen that \(B/a_{0}B\) is an A-module of finite length; so is B/b which is a quotient module of it. Further the B-module b contains, as a submodule, \(a_{0}B\) which is finitely generated; as \(b/a_{0}B\) is of finite length (as a submodule of \(B/a_{0}B\) ) and hence finitely generated, b is certainly a finitely generated B-module.

The above shows first that B is Noetherian. On the other hand, if p is a nonzero prime ideal of B, the ring B/p is an integral domain and of finite length and hence is a field (Algebra, Chapter VIII, § 6, no. 4, Proposition 9), so that p is maximal.

COROLLARY 1. For every prime ideal p of A, the set of prime ideals of B lying above p is finite.

Suppose first that \(\mathfrak{p} = (0)\) ; then the only prime ideal \(\mathfrak{q}\) of \(\mathbf{B}\) such that \(\mathfrak{q} \cap \mathbf{A} = (0)\) is (0); otherwise, writing \(\mathbf{S} = \mathbf{A} - \{0\}\) , \(\mathbf{S}^{-1}\mathfrak{q}\) would be a non-zero prime ideal of \(\mathbf{S}^{-1}\mathbf{B}\) (Chapter II, § 2, no. 5, Proposition 11) and \(\mathbf{S}^{-1}\mathbf{B}\) is just the field of fractions of \(\mathbf{B}\) , for it is a subring of \(\mathbf{L}\) containing \(\mathbf{K}\) (Algebra, Chapter V, § 3, no. 2, Proposition 3); whence an absurd conclusion. If now \(\mathfrak{p} \neq (0)\) , it follows from Proposition 5 that \(\mathbf{B}/\mathfrak{p}\mathbf{B}\) is a finite-dimensional vector space over the field \(\mathbf{A}/\mathfrak{p}\) , hence an Artinian ring and therefore has only a finite number of prime ideals (Chapter IV, § 2, no. 5, Proposition 9), which proves that there is only a finite number of prime ideals of \(\mathbf{B}\) containing \(\mathfrak{p}\) .

COROLLARY 2. The integral closure of A in L is a Dedekind domain.

This integral closure is an integrally closed Noetherian domain all of whose non-zero prime ideals are maximal; it suffices therefore to apply Theorem 1 of no. 2.

In particular:

COROLLARY 3. The integral closure of a Dedekind domain in a finite extension of its field of fractions is a Dedekind domain.

PROPOSITION 6. Let A be a Dedekind domain, K its field of fractions, L a finite extension of K and B the integral closure of A in L. Let p be a non-zero prime ideal of A, v the corresponding essential valuation of K and

\[
\mathrm{B} p = \prod_ {i} p _ {i} ^ {e (i)}
\]

the decomposition of the ideal Bp as a product of prime ideals. Then:

\begin{enumerate}
\def\labelenumi{(\alph{enumi})}
\item
  the prime ideals of B lying above p are the \(p_{i}\) such that \(e(i) > 0\) ;
\item
  the valuations \(v_{i}\) on \(L\) corresponding to these ideals \(p_{i}\) are, up to equivalence, the valuations on \(L\) extending \(v\) ;
\item
  \([\mathbf{B} / \mathfrak{p}_i\colon \mathbf{A} / \mathfrak{p}] = f(v_i / v)\) ;
\item
  \(e_i = e(v_i / v)\) (cf.~Chapter VI, § 8, no. 1, Definitions 1 and 2).
\item
  To say that a prime ideal q of B lies above p amounts to saying that \(q \supset p\) , hence that \(q \supset Bp\) and that q contains one of the \(p_i\) such that \(e(i) > 0\) (Chapter II, §1, no. 1, Proposition 1).
\item
  This follows, taking account of (a), from § 1, no. 8, Corollary to Proposition 12.
\item
  The residue field of \(v\) is identified with \(A / p\) and that of \(v_i\) with \(B / p_i\) (§ 1, no. 4, Corollary 1 to Proposition 6).
\item
  Let \(a\) (resp. \(a_i\) ) be a uniformizer for \(v\) (resp. \(v_i\) ). Then
\end{enumerate}

\[
\begin{array}{r l} a \mathrm{B} _ {p _ {i}} = a \mathrm{A} _ {p} \mathrm{B} _ {p _ {i}} = p \mathrm{A} _ {p} \mathrm{B} _ {p _ {i}} = p \mathrm{B}. \mathrm{B} _ {p _ {i}} & = \left(\prod_ {j} p _ {j} ^ {e (j)}\right) \mathrm{B} _ {p _ {i}} = \prod_ {j} \left(p _ {j} \mathrm{B} _ {p _ {i}}\right) ^ {e (j)} \\ & = \left(p _ {i} \mathrm{B} _ {p _ {i}}\right) ^ {e (i)} = a _ {i} ^ {e (i)} \mathrm{B} _ {p _ {i}} \end{array}
\]

since \(\mathfrak{p}_j\mathrm{B}_{\mathfrak{p}_i} = \mathrm{B}_{\mathfrak{p}_i}\) for \(j\neq i\) ; whence (d), since \(e(v_{i} / v) = v_{i}(a)\)

